\section{Lean integration without a second, weaker theorem prover}
\status{Integration specification. The cited core names were inspected; the adapter
interfaces below are proposed and have not been compiled.}

\subsection{Begin at a stable boundary, then make selected paths incremental}

The first Lean implementation should invoke \Grind{} on ordinary Lean goals
and supply newly proved local facts explicitly. This is less efficient than
an ideal native integration, but it permits the structural planner and
certificate pipeline to be tested without depending immediately on every
internal representation of \Grind{}.

There are concrete lower-level entry points to investigate next. In the
inspected Lean 4.34.0 source, \code{Lean.Meta.Grind.Params} includes
\code{extraFacts}, and \code{assertExtra} infers the types of those proof
expressions and adds them to the goal's fact machinery. The source also
exposes \code{GrindM.run}, \code{mkGoalCore}, and a \code{Result} containing
\code{failure?}, diagnostics, and counters.\cite{grindsource}
These are implementation hooks, not a promise of a stable public integration
contract.

In particular, feeding \code{Params.extraFacts} into a fresh run is not the
same as resuming an earlier saturation state. A true incremental adapter must
preserve the goal state and all relevant solver/context state, and must use
the normal internalization and propagation path when adding facts. Otherwise
an equality might exist in one table while being invisible to another solver.

A proposed versioned adapter has the following surface:
\begin{lstlisting}
GrindAdapter.open(snapshot, options) -> session
GrindAdapter.seed(session, checked_proofs) -> session
GrindAdapter.advance(session, work_quota) -> local_result
GrindAdapter.exportRelevant(session, demand) -> checked_summary
GrindAdapter.snapshot(session) -> resumable_state
GrindAdapter.restore(resumable_state) -> session
\end{lstlisting}

Keep all version-specific internals in one module. Unit-test seeded facts,
case-split rollback, proof extraction, metavariable changes, and equality
propagation through every enabled local theory. Until that adapter is reliable,
restarting a local tactic with explicit checked facts is preferable to a
fast but incoherent shared state.

\subsection{Use a typed semantic boundary for each external worker}

A reifier does more than turn Lean expressions into strings. For a polynomial
fragment, it produces a variable table, an exact polynomial, and a proof that
the polynomial's denotation under that table is the source expression. The
source type, operations, and their instances are part of the reification.
Unsupported subexpressions may become opaque atoms when that abstraction is
sound for the requested procedure.

A schematic contract is
\[
\operatorname{reify}(e)
  =(\text{environment }\rho,\text{ IR }t,
    \pi: e=\operatorname{eval}_{\rho}(t)).
\]
For an inequality or formula, the contract is an equivalence or a suitable
one-way implication between the source and reified statements. The direction
must match how the solver result will be used.

If congruence closure identifies two source terms, a worker can share one
polynomial variable only after preserving the equality proof needed to
justify that identification. It must not merge different variables because
their printed names happen to agree. Conversely, treating equal terms as
different opaque variables is usually a loss of proving power rather than
an unsound strengthening; this explains why the shared-state integration
matters.

For dependent expressions, equality classes must be homogeneous up to
justified type compatibility. Heterogeneous equality is not an excuse to
ignore type indices. Keep \code{HEq} reasoning and transport explicit, and
use Lean elaboration to validate casts. An initial external term language
should remain first-order and homogeneous instead of implementing an ad hoc
untyped approximation to all of Lean.

\subsection{Three levels of proof reconstruction}

The simplest backend returns ordinary Lean theorem applications and tactic
steps. This is appropriate for synthesized witnesses, explicit induction
motives, small nonnegative linear combinations, and recurrence identities.
The search can be complicated while replay remains short.

For larger algebraic results, return a typed certificate to a reflected Lean
checker with a theorem of the form
\[
\operatorname{check}(P,C)=\texttt{true}
\quad\Longrightarrow\quad \operatorname{Denotes}(P).
\]
The problem $P$ comes from the source reifier. The external solver is allowed
to choose only $C$, not silently replace $P$. In a conservative configuration,
reduce the checker inside the trusted proof-checking path rather than accept
an unchecked FFI result.

For general external ATP/SMT proofs, use a proof importer with explicit
supported rules. \code{lean-smt} is a concrete candidate: its documented
behavior includes returning proof-reconstruction holes as goals, and its
supported theories are not equivalent to the entire capability set of cvc5.
\cite{leansmt} Those holes enter the obligation graph. An external
\code{unsat} response does not bypass them.

\subsection{Keep the efficient existing bit-vector path}

A mixed goal involving opaque functions and bit-vectors should use the
existing shared-state \code{grind => bv_decide} route when applicable.
\Forge{} can schedule that route after structural steps or new lemmas have
made relevant facts available, but should not duplicate its encoding just to
fit a uniform architecture.\cite{release434}

A modular integer-to-bit-vector connection needs the correct modular
semantics. An integer embedding into a field can be injective; a fixed-width
bit-vector projection is not. Mixing those relationships in a generic
``numeric cast'' cache would be a serious error. Reifiers and proof bridges
must be indexed by the actual source and target operations.

\section{Higher-order structure and external reasoning}

\subsection{A focused source-term worker}

Before translating a goal to a first-order backend, expose its constructive
structure. For a function equality, attempt the relevant extensionality
principle and introduce a fresh argument. For a universally quantified goal,
introduce its binder. For a conjunction or constructor target, create the
corresponding subgoals. For an existential, construct a scoped witness and
then prove its specification.

A bounded type-directed term synthesizer can handle witness goals beyond the
affine fragment. Its grammar includes local terms, lambda abstraction,
application of retrieved lemmas/functions, constructors, and selected
recursor/fold schemas. Each candidate is elaborated before semantic checking.
Higher-order pattern unification handles tractable application patterns;
unrestricted higher-order unification is not assumed to be a cheap or complete
subroutine.

For example, to synthesize a function $h:\operatorname{List}(A)\to
\operatorname{List}(B)$ with the equations
\[
h([])=[],\qquad h(x::xs)=f(x)::h(xs),
\]
a list-fold schema already fixes the recursion structure. Search only needs
the empty case and the constructor step. A library index may retrieve map;
otherwise a recursor term provides an equivalent witness whose equations are
proved by reduction. This is a specific proposed worker, not something the
Python first-order list module already implements.

The architecture leaves room for more advanced type-inhabitation search.
Canonical is relevant prior work on proof synthesis involving higher-order
and dependent types.\cite{canonical} Such a worker should be integrated by
its proof-term interface, not by flattening away the dependencies that make
its synthesis task meaningful.

\subsection{First-order ATP workers are complementary}

After structural exposure, a saturation-based ATP may find useful theorem
combinations that a local tactic misses. Duper is an in-Lean proof-producing
option.\cite{duper} An external solver can also be used through a deliberate
translation and reconstruction layer. \code{lean-auto} provides relevant
translation infrastructure, including monomorphization support used by other
Lean automation projects.\cite{leanauto,leansmt}

The controller should not indiscriminately send every declaration in the
local environment. Select the dependency-relevant facts and instantiate
polymorphic schemas only at the types currently required by the proof demand.
Record which source theorem and universe/type instantiation each exported
formula came from, so reconstruction can apply the actual theorem.

If an abstraction intentionally forgets structure, document its direction:
every Lean countermodel should map to an abstract model when an unsatisfiability
argument relies on that property. A spurious abstract model then yields no
Lean refutation. The final protection remains proof reconstruction in the
original context, not an informal assertion that the translation is
``basically sound.''

\subsection{Candidate generation can be untrusted without being unconstrained}

An optional model, learned ranker, or language model can propose a lemma,
induction variable, proof sketch, or term. It cannot change the accepted
logical rules. It also cannot submit arbitrary Lean commands that introduce
new axioms or alter the environment under the guise of a proof term.
Restrict the import channel to a typed certificate schema or a term elaborated
inside a fixed environment. Resource and size limits protect the checker from
pathological proposals.

The design does not require an online model service. The executed prototypes
use deterministic enumeration, exact algebra, and local LP proposals. Their
algorithms can be useful even in a fully offline implementation.

\section{Soundness invariants and adversarial failure modes}

\subsection{The trust boundary}

For the proposed Lean tactic, correctness ultimately rests on the accepted
Lean environment and kernel checking of the completed proof. External search
code, optimization solvers, heuristics, and speculative graphs need not be
trusted to return correct answers. The implemented Python checkers are a
research prototype of this separation; their correctness is not established
by Lean's kernel.

An accepted source theorem must have no unresolved expression or universe
metavariables, no \code{sorry} dependency, and no unexpected new axioms relative
to the approved environment. Existing standard axioms and user-approved
assumptions are not silently removed or reclassified. A replay artifact should
make its dependencies inspectable, for example by normal Lean axiom reporting.

\subsection{The invariants to assert in code}

\begin{invariant}[Context validity]
Every proved fact is accompanied by a proof whose free variables and local
instances belong to its recorded context, or by an explicit abstraction and
instantiation that transports it into the consuming context.
\end{invariant}

\begin{invariant}[Candidate quarantine]
No candidate theorem, model equation, numerical conjecture, or unsolved proof
hole is admitted to the proved-fact state. Exploratory assumptions never serve
as proof evidence.
\end{invariant}

\begin{invariant}[Induction discipline]
An induction hypothesis is used only with the scope and smaller-argument
relationship supplied by its recursor. Generalized parameters may vary;
the structural induction index may not be replaced arbitrarily.
\end{invariant}

\begin{invariant}[Certificate binding]
A certificate is checked against the source-derived problem, with the same
variable mapping, types, assumptions, and target. A certificate for another
polynomial, another box, or a different quantifier scope is not interchangeable.
\end{invariant}

\begin{invariant}[Fail-closed arithmetic]
Numerical tolerances may guide search. Accepted equalities, coefficient signs,
partition boundaries, and proof-rule premises are checked exactly.
\end{invariant}

\subsection{An implementation review checklist}

The most dangerous implementation errors are not merely weak heuristics.
They are mistakes that accidentally turn search evidence into a stronger
assumption. Review at least the following cases explicitly:

\begin{longtable}{@{}>{\raggedright\arraybackslash}p{0.29\textwidth}>{\raggedright\arraybackslash}p{0.65\textwidth}@{}}
\toprule
Failure mode & Required defense \\
\midrule
\endhead
Self-justifying lemma & Prove candidates in a separate obligation; forbid dependency cycles unless justified by an explicit joint induction principle. \\
Induction hypothesis applied to the wrong list & Keep the tail/index rigid; let only genuinely generalized binders instantiate. \\
Branch fact leakage & Record split assumptions and scope; combine branches with a proof, not a global cache insertion. \\
Floating-point near-equality & Recompute exact rational identities; a small residual is a rejection, not a proof. \\
Wrong root problem & Bind the certificate to the original reified expression and assumptions, not an oracle-provided replacement. \\
Missing subdivision region & Derive child boxes from a full split tree and check both children. \\
Erased dependent type or cast & Preserve typed denotation proofs and explicit transport obligations. \\
Numerical field proof used for integer division & Use a source-semantics bridge or reject the fragment. \\
Opaque model element used as a witness & Reify an actual source term with the correct scope; otherwise reject the proposal. \\
Incomplete external reconstruction & Keep every hole as an unsolved obligation; no success until all are checked. \\
Stale state after metavariable assignment & Invalidate affected caches or restore a consistent snapshot containing the old identities. \\
\bottomrule
\end{longtable}

Only a subset of these has an executable analogue in the current miniature
languages. The tests cover rigid induction tails, unknown/cyclic rewrite
rules, malformed paths, sort mismatches, exact arithmetic corruption, wrong
constraints, witness corruption, and partition corruption. They do not test
a Lean-dependent-type reifier or a real branch-state adapter, because those
components have not been implemented.

\subsection{Proof size is a separate resource}

A solver may find a mathematical argument quickly yet produce a proof too
large to elaborate or check efficiently. Track the sizes of certificates,
polynomial supports, rewrite traces, and final Lean proof terms separately
from search time. Share proof subterms, and prefer a small checked reflection
lemma to expanding an enormous algebraic identity into individual rewrites.

A proof DAG can be topologically emitted as local lemmas. Minimize a replay
script only by rerunning and rechecking the simplified version. Never delete
a dependency because it appears unused in a heuristic trace while remaining
inside the actual proof term.

\section{A concrete implementation plan}

\subsection{Module boundaries}

A Lean package could use the following organization. These are proposed file
names, not claims about files already present in this archive.
\begin{lstlisting}
Forge/Core/GoalSnapshot.lean
Forge/Core/Obligation.lean
Forge/Core/CheckedFact.lean
Forge/Core/Replay.lean
Forge/Search/Controller.lean
Forge/Search/RuleIndex.lean
Forge/Search/Structural.lean
Forge/Induction/CallAnalysis.lean
Forge/Induction/Generalize.lean
Forge/Induction/LemmaFrontier.lean
Forge/Synthesis/AffineWitness.lean
Forge/Synthesis/Recurrence.lean
Forge/Instantiation/ConstrainedMatch.lean
Forge/Instantiation/ModelGuided.lean
Forge/Reify/Polynomial.lean
Forge/Certificate/Farkas.lean
Forge/Certificate/Bernstein.lean
Forge/Backend/Grind.lean
Forge/Backend/SOS.lean
Forge/Backend/SMT.lean
Forge/Tactic.lean
\end{lstlisting}

The core modules depend only on Lean and a small shared representation.
Mathlib-specific reifiers and tactics belong behind separate imports.
External processes consume a versioned certificate/problem format and return
only untrusted proposals. The package remains useful when those processes
are unavailable.

\subsection{Milestones with objective exit criteria}

\textbf{First milestone: proof-state controller and baseline parity.}
Implement exact-state restore, branch-scoped obligations, replay records, and
a \Grind{} tactic adapter. Its restricted mode must reproduce a pinned baseline
suite without introducing new axioms or losing source assumptions.

\textbf{Second milestone: structural induction and accumulator generalization.}
Port the miniature list algorithm using real Lean expressions, then extend it
to a small set of inductive types. The decisive tests are not just successes:
failed generalizations, dependent hypotheses, and wrong-tail induction
hypotheses must be rejected or left unsolved.

\textbf{Third milestone: exact certificates.}
Implement the affine/Farkas certificate theorem and the recurrence certificate
reconstruction. Integrate the existing SOS backend behind the same checked-fact
interface. Add Bernstein checking only after the source box and polynomial
reifiers have their soundness contracts.

\textbf{Fourth milestone: missing terms and missing lemmas.}
Add arithmetic constrained matching, typed lemma-frontier abstraction, and
counterexample filtering. Evaluate generalization separately from helper-lemma
discovery so the contribution of each is visible.

\textbf{Fifth milestone: incremental internals and optional external workers.}
Introduce a version-pinned \Grind{} session adapter only where measurements
show repeated setup costs dominate. Add SMT and higher-order workers without
expanding the trusted base or pretending every external proof rule is already
supported.

\textbf{Sixth milestone: hardening and competitive evaluation.}
Use holdout benchmarks, per-worker ablations, failure classifications,
certificate mutation tests, proof-size measurements, and clean-build replay.
A release claiming significant superiority should be based on these results,
not on the manufactured successes in this article.

\subsection{User interface and explanations}

A possible syntax, explicitly \emph{proposed}, is:
\begin{lstlisting}
-- Proposed syntax, not supplied as an executable tactic:
forge [relevant_definitions, useful_theorems]
forge? (induction := true) (witnesses := .affine)
forge (external := false) (maxLemmaDepth := 2)
\end{lstlisting}

\code{forge?} should emit a stable proof script or certificate replay plus any
local helper lemmas it actually needed. A failure report should distinguish
``no induction motive found,'' ``candidate lemma refuted by an executable
example,'' ``nonlinear certificate search exhausted,'' ``unsupported source
operation,'' and ``external proof reconstruction incomplete.'' It should
identify the unsolved obligation rather than merely advise a larger timeout.

The replay should not rerun stochastic discovery or require the original
external solver. A proof that only succeeds while a particular model service
is available is an unsatisfactory end product when an explicit certificate
can be retained instead.
